\documentclass[11pt]{article}

% Literature review notes for the Related Work section of the TraceAssist
% paper. Standalone; builds with plain pdflatex.
\usepackage[margin=1in]{geometry}
\usepackage[T1]{fontenc}
\usepackage{xspace}
\usepackage{url}
\usepackage[hidelinks]{hyperref}
\usepackage{enumitem}
\setlist{nosep,leftmargin=*}

\newcommand{\sysname}{TraceAssist\xspace}

\title{Literature Review: Work Related to \sysname}
\author{}
\date{October 2026}

\begin{document}
\maketitle

\noindent Three papers are closest to \sysname. Each one combines an LLM with
evidence prepared outside the LLM, as \sysname does, but they differ in what
the evidence is and what the LLM is asked to do. Details below come from each
paper's own text (TAAF from the local PDF, \texttt{RAG\_Qs/2601.02632v1.pdf};
the other two from their arXiv abstracts).

\section{TAAF: A Trace Abstraction and Analysis Framework Synergizing
Knowledge Graphs and LLMs}

\textbf{Citation.} A.~Ezaz, G.~Khodabandeh, M.~Babaei, N.~Ezzati-Jivan.
ICSE 2026. arXiv:2601.02632, DOI 10.1145/3744916.3787832.
\url{https://arxiv.org/abs/2601.02632}

\textbf{Summary.} TAAF answers natural-language questions about Linux kernel
traces. It builds a time-indexed Trace Compass \emph{state system} from the
trace, and for each question extracts a query-specific \emph{knowledge graph}
of threads, CPUs and resources; an LLM reasons over that graph to answer. The
authors introduce \textbf{TraceQA-100}, 100 expert-written questions on LTTng
traces of the SciMark 2.0 Java benchmark (about 34 million kernel events).
Each question is posed on nine slices: 1, 10 and 100\,s windows at the start,
middle and end of the trace. Questions are tagged single-hop or multi-hop.
They compare three inputs: raw events (qualitative only, since they exceed
the context window), the state system's numeric output (the baseline), and
the knowledge graph (TAAF), with GPT-4.1 nano, GPT-4o and o4-mini, plus
Gemini 2.5 and Claude models as a cross-vendor probe. TAAF improves accuracy
by up to 31.2\%, mostly on multi-hop and causal questions.

\textbf{Similar to \sysname.}
\begin{itemize}
  \item Same problem: natural-language questions about LTTng kernel traces,
        motivated by traces being too large for an LLM's context.
  \item Same foundation: Trace Compass's state system.
  \item Same finding: multi-hop questions are harder, and structured
        grounding helps most there. Our hub-scope labels follow TAAF.
\end{itemize}

\textbf{Different from \sysname.}
\begin{itemize}
  \item TAAF gives the LLM a graph and leaves the reasoning and arithmetic to
        it; \sysname's analyzers compute the answer's measurements (wait
        chains, per-round finish times, lock holds) and the LLM explains them.
  \item TAAF answers over slices of at most 100\,s; \sysname analyses the
        whole trace (up to 30 minutes, 87 million events).
  \item TAAF has no documentation retrieval; \sysname adds RAG.
  \item TAAF's questions describe scheduling behaviour of a benchmark;
        ours diagnose an injected performance bug, including root cause and
        remediation.
\end{itemize}

\textbf{Use in Related Work.} The closest prior work; position \sysname as
moving from a graph the LLM must reason over to computed findings the LLM
only explains, and from short slices to whole traces.

\section{RCACopilot: Automatic Root Cause Analysis via Large Language Models
for Cloud Incidents}

\textbf{Citation.} Y.~Chen, H.~Xie, M.~Ma, Y.~Kang, X.~Gao, L.~Shi, Y.~Cao,
X.~Gao, H.~Fan, M.~Wen, J.~Zeng, S.~Ghosh, X.~Zhang, C.~Zhang, Q.~Lin,
S.~Rajmohan, D.~Zhang, T.~Xu. EuroSys 2024. arXiv:2305.15778.
\url{https://arxiv.org/abs/2305.15778}

\textbf{Summary.} RCACopilot is an on-call system for cloud incidents at
Microsoft. Each incident is matched by its alert type to an \emph{incident
handler}, which collects the relevant runtime diagnostic information; an LLM
then predicts the root-cause category and writes an explanation. On one
year of Microsoft incidents it reaches RCA accuracy of up to 0.766, and its
diagnostic-collection component has been in production for over four years.

\textbf{Similar to \sysname.}
\begin{itemize}
  \item The same division of labour: deterministic components gather the
        evidence (handlers; for us, the planner and analyzers), and the LLM
        turns it into an explanation.
  \item Both target root cause analysis for engineers who would otherwise
        inspect logs and traces by hand.
\end{itemize}

\textbf{Different from \sysname.}
\begin{itemize}
  \item Cloud incidents and alerts, not kernel traces; the evidence is
        diagnostic data selected by alert type, not measurements computed
        from a trace.
  \item The LLM classifies incidents into root-cause categories; \sysname
        answers open questions, including quantitative and remediation
        questions.
\end{itemize}

\textbf{Use in Related Work.} Shows that the ``deterministic evidence, LLM
explanation'' design works in production for incidents; \sysname applies it
to kernel traces, where the evidence must itself be computed.

\section{Opal: A Modular Framework for Optimizing Performance using
Analytics and LLMs}

\textbf{Citation.} M.~Zaeed, T.~Z.~Islam, V.~In\dj{}i\'c. arXiv:2510.00932,
October 2025 (preprint).
\url{https://arxiv.org/abs/2510.00932}

\textbf{Summary.} Opal feeds performance analytics into LLM prompts so the
model can suggest code optimizations for GPU kernels. Its insights come from
hardware counters and Roofline analysis, linked to stall events. Across
1{,}640 experiments on real GPU kernels, even a single insight source gave a
speedup in over 98.5\% of cases, with average speedups of 19.34--52.3\%; the
generated code was correct in all but one case, a failure caused by a vague
diagnostic.

\textbf{Similar to \sysname.}
\begin{itemize}
  \item Grounds the LLM in performance evidence computed by analysis tools
        rather than raw data.
  \item Targets the step that remains manual and expert-driven: turning
        diagnostic output into an explanation or a fix (our remediation
        questions).
  \item Its one failure, a vague diagnostic leading to an unsafe
        suggestion, supports our design of precise findings with confidence
        levels and caveats.
\end{itemize}

\textbf{Different from \sysname.}
\begin{itemize}
  \item Hardware-counter profiles of GPU kernels, not kernel traces of
        multithreaded programs; it does not reason about waits between
        threads.
  \item The output is optimized code, evaluated by speedup; \sysname's
        output is answers to questions, evaluated by correctness.
\end{itemize}

\textbf{Use in Related Work.} Evidence from the performance-engineering side
that analytics-grounded prompts make LLM suggestions reliable.

\section{Also considered}

\begin{itemize}
  \item \emph{Stalled, Biased, and Confused: Uncovering Reasoning Failures in
        LLMs for Cloud-Based Root Cause Analysis} (Riddell et al., FORGE
        2026, arXiv:2601.22208). Studies six LLMs' reasoning failures in
        cloud RCA, including sensitivity to input data modalities. A
        search snippet said removing traces sometimes improved accuracy,
        which would support our raw-trace baseline results, but the
        abstract does not state it; check the full paper before citing.
  \item \emph{Exploring LLM-based Agents for Root Cause Analysis}
        (arXiv:2403.04123): a ReAct agent with tools for incident RCA. Not
        read in detail.
\end{itemize}

\section{Suggested positioning}

Prior work grounds LLMs in structure (TAAF's knowledge graphs), in
diagnostic data selected per incident (RCACopilot), or in profiling
analytics (Opal). \sysname differs by computing the answer's measurements
from the whole kernel trace before the LLM sees anything, and by combining
those findings with retrieved documentation.

\end{document}
